% Arquivo gerado automaticamente: esqueleto LaTeX do artigo
\documentclass[a4paper,12pt]{article}
\usepackage[T1]{fontenc}
\usepackage[utf8]{inputenc}
\usepackage[portuguese]{babel}
\usepackage{csquotes}
\usepackage{microtype}
\usepackage{hyperref}

\title{Título provisório: Índice Sintético de Viabilidade Territorial (ISVT) e o Turismo Rural}
\author{Seu Nome Aqui \\ \small{Instituição / Programa}}
\date{\today}

\begin{document}
\maketitle

\begin{abstract}
Este é um rascunho inicial. O trabalho propõe o Índice Sintético de Viabilidade Territorial (ISVT) para analisar o potencial de desenvolvimento do ecoturismo de base comunitária na Bahia e no Rio Grande do Sul. Substitua este resumo por uma versão finalizada.
\end{abstract}

\section{Introdução}

O turismo consolida-se como um vetor estratégico de desenvolvimento socioeconômico e de descentralização da renda em países em desenvolvimento, especialmente em regiões marcadas pela retração de atividades industriais e agrícolas tradicionais. Nesse cenário, a transição para modalidades alternativas — com destaque para o turismo rural, o ecoturismo, o agroturismo e o turismo de base comunitária — surge como uma via para a diversificação da matriz produtiva. Essas atividades promovem a valorização cultural local, a permanência das famílias no campo, a inclusão social e o fortalecimento do orgulho comunitário. Diferentemente do turismo de massa, tais vertentes valorizam o "saber fazer" local, a autenticidade das paisagens e a herança cultural intangível de cada território, gerando ocupações complementares à agricultura familiar sem descaracterizar a identidade camponesa.

Ressalta-se, contudo, que a mera existência de atrativos naturais ou culturais não garante o desenvolvimento regional (Caetano, Wegner \& Xavier, 2023). Atualmente, a base econômica de muitas comunidades é fragilizada pela falta de planejamento público e pela escassez de investimentos. Esses recursos são fundamentais para fomentar um setor que, apenas nas zonas rurais, tem o potencial de impactar 25 milhões de brasileiros (IBGE, 2022).

Sustenta-se que, apesar desse expressivo potencial, a expansão do turismo rural no Brasil ocorre sob uma profunda heterogeneidade estrutural. Esse panorama revela a coexistência de setores modernos e altamente produtivos ao lado de bolsões de subsistência. Essa disparidade reflete-se diretamente na distribuição espacial dos recursos de fomento e crédito no território: enquanto municípios-polo e capitais regionais concentram a maior parte dos investimentos privados e da infraestrutura de grande escala, pequenas localidades vocacionadas ao ecoturismo comunitário enfrentam severas restrições de liquidez financeira.

Diante da escassez de investimentos e da ausência de políticas integradas para mitigar tais assimetrias, o sucesso operacional de um destino turístico depende do nível de articulação organizacional e da capacidade do território de converter seus ativos em produtos comercializáveis e sustentáveis. Essa conversão é crucial para que a diversificação turística gere valor cultural, preserve a biodiversidade e promova o avanço socioeconômico local. Afinal, as limitações no acesso a mecanismos formais de financiamento e a falta de infraestrutura de suporte (como telecomunicações e transporte) constituem, hoje, os dois maiores gargalos para o desenvolvimento turístico sustentável em áreas não urbanas (UN Tourism, 2023).

O objetivo central deste trabalho é analisar a estrutura dos municípios para o desenvolvimento do ecoturismo de base comunitária, avaliando em que medida a circulação de capital financeiro no campo dialoga com a oferta produtiva e a atratividade dos destinos. A escolha dos estados justifica-se por seus perfis geográficos e institucionais marcadamente contrastantes. O Rio Grande do Sul apresenta um histórico consolidado de rotas rurais e cooperativismo de crédito integrado; por outro lado, a Bahia conta com polos turísticos de grande escala concentrados na faixa litorânea, enfrentando o desafio de interiorizar as receitas do setor em direção ao semiárido e à Chapada Diamantina.

A presente pesquisa propõe, portanto, a formulação do Índice Sintético de Viabilidade Territorial (ISVT) como uma ferramenta de geointeligência aplicada à governança do turismo na Bahia (BA) e no Rio Grande do Sul (RS). O ISVT pretende mapear esses contrastes latentes para subsidiar o planejamento de políticas públicas integradas de fomento e desenvolvimento regional.

\section{Metodologia (esboço)}
Descreva aqui as fontes de dados (censo, cadastros, séries administrativas), a construção dos indicadores, a normalização e a agregação que comporão o ISVT. Sugiro incluir um fluxograma e uma tabela resumo com os indicadores por dimensão.

\section{Resultados esperados}
Apresente o que se espera identificar com o ISVT: zonas com maior viabilidade, gargalos de financiamento, correlação entre atratividade e acesso a crédito, etc.

\section{Conclusão (provisória)}
Insira uma conclusão provisória com recomendações de política e limitações do estudo.

\bibliographystyle{plain}
\bibliography{references}

\end{document}
